
接着在
$\bigl([0,1]^{|\vec a|}\times
\mathcal M^{\mathrm{main},\vec a}_{k+1}
(\beta;\vec P)\bigr)^{\mathfrak s}$
上定义等价关系 $\sim'$。

\begin{definition}{7.3.28}\label{def:7-3-28}
设
$(\vec t,\mathbf x)$、
$(\vec t',\mathbf x')$
属于
$[0,1]^{|\vec a|}\times
\mathcal M^{\mathrm{main},\vec a}_{k+1}(\beta)$。
若存在相容条件~\ref{cond:7-3-13}
所述的 $\vec a^1,\vec a^2$，
以及
\[
(\vec t_-,\mathbf x_-)\in
[0,1]^{|\vec a^2|}\times
\mathcal M^{\mathrm{main},\vec a^{2\prime}}_{
k+1-|\vec a^1|}(\beta),
\]
使得
\[
I^1_{\vec a^1,\vec a^2}(\vec t_-)
=\vec t=\vec t',
\qquad
\operatorname{forget}_{\vec a^1,\vec a^2}
(\mathbf x)
=\operatorname{forget}_{\vec a^1,\vec a^2}
(\mathbf x')
=\mathbf x_-,
\]
则称二者等价，记作
$(\vec t,\mathbf x)\sim(\vec t',\mathbf x')$。
\translatornote{源书把 $\mathbf x_-$ 所在模空间
写成
$\mathcal M^{\mathrm{main},\vec a}_{
k+1-|\vec a^2|}(\beta)$；
按相容条件 7.3.13 前的遗忘映射靶空间，
应删去 $|\vec a^1|$ 个标记点，
并把所余序列记作 $\vec a^{2\prime}$。
见疑误 A-0548。}
\end{definition}

这一等价关系在
$[0,1]^{|\vec a|}\times
\mathcal M^{\mathrm{main},\vec a}_{k+1}
(\beta;\vec P)$
的仓西邻域上诱导等价关系。
再由相容条件~\ref{cond:7-3-13}，
它在零集
$\bigl([0,1]^{|\vec a|}\times
\mathcal M^{\mathrm{main},\vec a}_{k+1}
(\beta;\vec P)\bigr)^{\mathfrak s}$
上诱导等价关系 $\sim'$。
记其等价类集合为
\begin{equation}
\bigl(
([0,1]^{|\vec a|}\times
\mathcal M^{\mathrm{main},\vec a}_{k+1}
(\beta;\vec P))^{\mathfrak s}
\bigr)_{\mathrm{Coll}}.
\tag{7.3.29}\label{eq:7-3-29}
\end{equation}
由引理~\ref{lem:7-3-27}，
式~\eqref{eq:7-3-29} 有三角剖分。
还有分段光滑的评价映射
\[
\operatorname{ev}\colon
\bigl(
([0,1]^{|\vec a|}\times
\mathcal M^{\mathrm{main},\vec a}_{k+1}
(\beta;\vec P))^{\mathfrak s}
\bigr)_{\mathrm{Coll}}
\longrightarrow L^{k+1-|\vec a|}.
\]

以式~\eqref{eq:7-3-29} 的商空间
代替原来的零集，
便得到与命题~\ref{prop:7-3-10} 相同的结论。
相容条件~\ref{cond:7-3-13}
不再作为额外条件出现，
因为它已由这一商空间的良定义性取代。
实际用来定义 $\mathfrak h$ 的
正是此商空间。
由引理~\ref{lem:7-3-22} 第~(3) 项，
情形 1 对应的边界贡献
于是严格等于零。

最后讨论
$\mathcal M^{\mathrm{main}}_{
k+1-|\vec a|}(\beta)$
为空的情形，
即相容条件~\ref{cond:7-3-14}
所涉及的情形。
此时考虑第 $0$ 个标记点处的评价映射
\begin{equation}
\operatorname{ev}_0\colon
\mathcal M^{\mathrm{main}}_{k+1}(\beta_0)
\longrightarrow L.
\tag{7.3.30}\label{eq:7-3-30}
\end{equation}
并用下列交换图代替图式~\ref{diag:7-3-1}：

\par\medskip\noindent\begin{minipage}{\textwidth}
\centering
\resizebox{\textwidth}{!}{$\displaystyle
\begin{array}{ccc}
\mathcal M^{\mathrm{main},\vec a^1}_{k_1+1}(\beta_0)
{}_{\operatorname{ev}_0}\times_{\operatorname{ev}_i}
\mathcal M^{\mathrm{main},\vec a^2}_{k_2+1}(\beta_0)
&\longrightarrow&
\mathcal M^{\mathrm{main},\vec a}_{k+1}(\beta_0)
\\[0.4em]
{\scriptstyle\operatorname{ev}_0}\Big\downarrow
&&
\Big\downarrow{\scriptstyle\operatorname{ev}_0}
\\[0.4em]
L&=&L
\end{array}$}
\par\smallskip
{\renewcommand{\thefigure}{7.3.3}\refstepcounter{figure}
\normalfont\bfseries 图式~\thefigure\label{diag:7-3-3}}
\end{minipage}\par\medskip

图式~\ref{diag:7-3-3} 的第一根竖箭头中，
$\operatorname{ev}_0$
作用于第二个因子
$\mathcal M^{\mathrm{main},\vec a^2}_{
k_2+1}(\beta_0)$。
随后以图式~\ref{diag:7-3-3}
代替图式~\ref{diag:7-3-1}，
得到与图式~\ref{diag:7-3-2} 类似的图式。
这样便找到扰动 $\mathfrak s$
及其零集
$\mathcal M^{\mathrm{main},\vec a}_{k+1}
(\beta_0;\vec P^+)^{\mathfrak s}$。
再仿照定义~\ref{def:7-3-28}
在该零集上定义等价关系。

注意，除 $k=2$ 外，
式~\eqref{eq:7-3-30} 的纤维为正维
\translatornote{源书此处误引“(7.3.27) 的纤维”；
7.3.27 是引理，紧前的 (7.3.30)
才是评价映射。见疑误 A-0549。}。
因此可按相同方式讨论，
完成定理~\ref{thm:7-3-1} 的证明。
定理~\ref{thm:7-3-1} 证毕。\hfill$\square$\par

同理还可证明：
第 4 章和第~\ref{sec:7-2} 节构造的
各类滤过 $A_\infty$ 同态及其间同伦，
都是同伦含幺的；
第 5 章构造的滤过 $A_\infty$ 双模
及其间同态，也都是同伦含幺的。
